\section{Introduction}
A finite statistical memory is a collection of possible continuations,
not a collection of numerical coordinates. A low-dimensional space of
predictions can require many positive generators. Conversely, insisting
that every internal state be an observed residual can exclude much smaller
stochastic realizations. There is a third constraint: positive
factorizations at successive cuts must arise from the same causal machine.
The distinctions between dimension, positivity, and compatibility are the
subject of this paper.

We work with complete finite response arrays, known primitive probabilities,
and an explicit list of charged cuts. Arbitrary stochastic rows are atomic
operations. Every retained selector, observation buffer, and data-dependent
protocol state is included in the register. The program and the declared
clock are fixed before observations. These conventions distinguish an exact
state-count theorem from a finite-random-bit implementation or an acquired
calibration theorem.

\subsection{The realization question}
At a cut $t$, normalized residual continuations form a finite set
$\cR_t$ in a causal polytope $\cP_t$. Intersections of their minimal faces
partition them into support components. If component $j$ has affine
dimension $r_{t,j}-1$, every stochastic implementation needs at least
$s_t=\sum_jr_{t,j}$ states. Residual vertices attain this bound when each
component hull is a simplex. That sufficient condition does not describe
all rank-saturating realizations.

Theorem~\ref{thm:coherence} gives a necessary and sufficient condition for
attaining the whole profile $s_t$. Each component must admit an enclosing
simplex with $r_{t,j}$ normalized causal vertices in its affine slice. The
vertices need not be residuals. Their one-step restrictions must satisfy
common nonnegative shift equations with the next cut's vertices. Thus the
condition separates static enclosure from simultaneous causal realization.
For explicit algebraic arrays it becomes a finite existential system over
the real numbers. This is an exact equality criterion, not a polynomial-time
algorithm for a succinct controller or a solution of autonomous
shared-row minimization.

\subsection{The growing experiment}
An input $x\in\{-1,1\}^n$ is acquired one coordinate at a time. After all
coordinates have been acquired, a query $j\in[n]$ is supplied. The required
binary response is
\begin{equation}\label{eq:query-intro}
 \Prb(B=b\mid x,j)=\frac{1+b\eta x_j}{2},\qquad b\in\{-1,1\}.
\end{equation}
The signal $\eta$ is specified, and the law is required for every input
and query. Write $K_n$ for the smallest state count at the query barrier,
$W_n^{\rm ad}$ for the smallest whole-schedule peak with adaptive
acquisition, and $W_n^{\rm fo}$ for the fixed-order peak. In the latter two
quantities all persistent choices during acquisition are charged. They
satisfy $K_n\le W_n^{\rm ad}\le W_n^{\rm fo}$.

Let $h_2$ be binary entropy in bits and put
\begin{equation}\label{eq:rate-intro}
 I(u)=1-h_2((1+u)/2),\qquad 0\le u\le1.
\end{equation}
The main quantitative conclusion, Theorem~\ref{thm:rate}, is
\begin{equation}\label{eq:headline-rate}
 \lim_{n\to\infty}\frac{\log_2K_n(\eta)}n
 =\lim_{n\to\infty}\frac{\log_2W_n^{\rm ad}(\eta)}n
 =\lim_{n\to\infty}\frac{\log_2W_n^{\rm fo}(\eta)}n=I(\eta)
 \quad(0<\eta<1).
\end{equation}
The checkpoint upper error is $O_\eta(\log n)$ bits. The streaming error
is $O_\eta(\sqrt{n\log n})$ bits. A fixed codebook is chosen in the
existence proof, so there is no persistent public random seed hidden from
the state count. For every rowwise total-variation tolerance $\varepsilon$,
Theorem~\ref{thm:approx-rate} gives the rate $I((\eta-2\varepsilon)_+)$.

The exponent in \eqref{eq:headline-rate} is not new as a random-access
coding exponent. Ambainis, Nayak, Ta-Shma and Vazirani
\cite{ANTVpre,ANTV} prove the matching first-order communication rate
for success probability at least $(1+\eta)/2$. We prove a version fixing
\emph{each} conditional response in \eqref{eq:query-intro}, and then
account for the largest retained register during a one-pass implementation.
A success inequality by itself is not an exact conditional channel.
Likewise, a short message by itself is not a space bound for computing
that message from a stream. The proofs below keep both distinctions.

Finite-dimensional information is also obtained. An elementary online
simplex construction gives the exact acquisition profile $t+1$ for
$0<\eta\le(2n-1)^{-1}$, improving the smaller signal regime of the
preceding development. Any $n+1$-state checkpoint factorization requires
$\eta\le1/n$. If a Hadamard matrix of order $n+1$ exists, a different
online encoder attains that entire interval with peak $n+1$
(Theorem~\ref{thm:hadamard}). In particular, the threshold is sharp on
the infinite sequence $n=2^k-1$. The latter construction has peak $n+1$;
it is not claimed to have only $t+1$ states at every earlier cut.

\Needspace{14\baselineskip}
\subsection{One notation and resource convention}
\begin{center}
\begin{tabular}{@{}p{.23\textwidth}p{.69\textwidth}@{}}
\toprule
Symbol & Meaning\\
\midrule
$C$, $R_h$, $\cP_t$ & Complete causal law, normalized residual, and ambient causal polytope.\\
$r_{t,j}$, $s_t$ & Rank of support component $j$, and their sum.\\
$B_{t,j,k}$ & An internal normalized continuation, not necessarily a residual.\\
$K_n$, $W_n^{\rm fo}$, $W_n^{\rm ad}$ & Query-cut, fixed-order peak, and adaptive-acquisition peak state counts.\\
$\eta$, $\varepsilon$, $I$ & Specified signal, uniform row-TV tolerance, and entropy rate in bits.\\
\bottomrule
\end{tabular}
\end{center}
State counts are cardinalities; their binary memory cost is the ceiling
of their base-two logarithm. The query is an atomic current input and is
answered without retaining an uncharged copy. Input-dependent scheduling
information has to be retained. A terminal output can use two labels;
this constant does not affect any rate. An externally scheduled overwrite
may draw fresh independent randomness, but past randomness survives only
through the retained label.

\subsection{Relationship with the preceding foundations development}
The causal-experiment, predictive-quotient, minimax-contact, and physical
comparison results of the preceding development \cite{GTF29} are preserved
in a separate complete mathematical companion. The present article proves
its new claims without asking the reader to reconstruct a proof from that
volume. Section~\ref{sec:literature} identifies the relations with positive
realization, residual automata, and random-access coding. The old
known-target physical memory theorems remain unchanged; neither a new
fully adaptive collision theorem nor closure of the historical analytic
pipeline is inferred from \eqref{eq:headline-rate}.
